% SOURCE_AUTHORITY: sga5_fr_workpass.tex, lines 13205--13229 (LF-normalized unit).
% SOURCE_SHA256: 791F4EFFC5E02832D5D77ED03518C8156D6F07E4C8238B03545DB93D883FBB28
Por último, para estabelecer (7.13) basta utilizar a fórmula de Euler--Poincaré (7.2) para o complexo de feixes $i_!(F)$:
\[
\EP_\Delta(V,F)
=
\clK{\Delta}(\RG_{!V}(F))
=
\clK{\Delta}(\RG_C(i_!(F)))
=
\EP(C)\clK{\Delta}(F_{\bar\eta})
-\sum_{x\in C(k)}\varepsilon_x^\Delta(i_!(F)).
\]
Ora,
\[
\sum_{x\in C(k)}\varepsilon_x^\Delta(i_!(F))
=
\sum_{x\in V(k)}\varepsilon_x^\Delta(F)
+\sum_{x\in Z}\varepsilon_x^\Delta(F).
\]
Mas, para $x\in Z$, tem-se $(i_!(F))_{\bar x}=0$ e, por conseguinte, para tal $x$ tem-se
\[
\varepsilon_x^\Delta(i_!(F))
=
\alpha_x^\Delta(F)+\clK{\Delta}(F_{\bar\eta}),
\]
de onde resulta a fórmula (7.13).
\end{prooftext}
